\chapter{Dihydrocodeine}
\index{Dihydrocodeine}

\section*{Definition and Overview}
Dihydrocodeine is a prescription opioid medicine used to relieve moderate to severe pain. It is related to codeine but is somewhat stronger, and it works by acting on opioid receptors in the brain and spinal cord to reduce the perception of pain. Dihydrocodeine is used for short-term pain relief after surgery or injury, and occasionally for other pain under specialist guidance. Like all opioids, it can cause dependence, tolerance, and overdose, and it is a controlled medicine in many countries. This chapter explains how dihydrocodeine works, how it is used safely, and its risks.

\section*{Epidemiology}
Opioid medicines, including dihydrocodeine, are widely prescribed for pain, although there is now greater caution about their use because of dependence and overdose risks. some countries have seen rising rates of opioid-related harm, prompting tighter regulation, smaller pack sizes, and prescribing restrictions. Dihydrocodeine is available in many countries but is less commonly used than other opioids. Its use is appropriate for short-term severe pain that has not responded to other medicines.

\section*{Aetiology and Pathophysiology}
\begin{itemize}
    \item \textbf{Mechanism:} dihydrocodeine binds to opioid receptors, blocking pain signals and altering the emotional response to pain
    \item \textbf{Onset and duration:} effects begin within 30--60 minutes and last several hours
    \item \textbf{Metabolism:} it is partly converted to an active form in the body
    \item \textbf{Tolerance:} repeated use leads to reduced effect, requiring higher doses
    \item \textbf{Dependence:} physical dependence develops with regular use, causing withdrawal on stopping
    \item \textbf{Addiction:} some people develop compulsive use despite harm
\end{itemize}

\section*{Clinical Features}
\begin{itemize}
    \item Used for moderate to severe pain that is not controlled by simpler pain relievers
    \item Common side effects: constipation, nausea, vomiting, drowsiness, and dizziness
    \item Effects on breathing, which can be dangerous at high doses or with other sedatives
    \item Reduced concentration and reaction time, affecting driving
    \item Withdrawal symptoms on stopping after regular use: anxiety, sweating, cramps, and diarrhoea
    \item Risk of overdose with excessive doses or combined with alcohol or other sedatives
\end{itemize}

\section*{Differential Diagnosis}
\begin{itemize}
    \item Dihydrocodeine versus other pain medicines, such as paracetamol, NSAIDs, and other opioids
    \item Short-term versus long-term pain management strategies
    \item Opioid therapy versus non-drug approaches, including physiotherapy and psychological treatment
\end{itemize}

\section*{When to Seek Medical Care}
Dihydrocodeine is a prescription medicine and should be used exactly as directed and for the shortest time needed. See a doctor if pain persists, if side effects are troublesome, or before stopping after regular use. Never take it with alcohol or other sedating medicines without advice.

\textbf{Contact your local emergency services for:}
\begin{itemize}
    \item Difficulty breathing, severe drowsiness, or unresponsiveness (possible overdose)
    \item Taking dihydrocodeine with alcohol or other sedatives with serious effects
    \item Severe confusion or collapse
    \item Any sign of an allergic reaction
\end{itemize}

\section*{Diagnosis and Investigations}
\begin{itemize}
    \item \textbf{Medical assessment:} review of the pain, its cause, and appropriate treatment
    \item \textbf{Review of medicines:} checking for interactions and other sedating drugs
    \item \textbf{Risk assessment:} history of dependence, substance use, or breathing problems
    \item \textbf{Liver and kidney function:} dose adjustment may be needed
    \item \textbf{Monitoring:} review of pain control, side effects, and signs of dependence
\end{itemize}

\section*{Management}
\begin{itemize}
    \item \textbf{Shortest duration:} used for days rather than weeks wherever possible
    \item \textbf{Lowest effective dose:} minimises side effects and dependence risk
    \item \textbf{Regular review:} reassessment of the need for ongoing opioids
    \item \textbf{Constipation management:} laxatives, fluids, and fibre are often needed
    \item \textbf{Avoiding alcohol:} alcohol greatly increases sedation and overdose risk
    \item \textbf{Gradual withdrawal:} tapering under medical supervision after regular use
    \item \textbf{Alternatives:} non-opioid medicines and non-drug pain management for long-term pain
\end{itemize}

\section*{Complications}
\begin{itemize}
    \item Dependence and addiction
    \item Overdose, including fatal respiratory depression
    \item Constipation and other gastrointestinal effects
    \item Drowsiness and falls, especially in older people
    \item Withdrawal symptoms on stopping
    \item Impaired driving and workplace safety
\end{itemize}

\section*{Prognosis and Outcomes}
Dihydrocodeine is effective for short-term severe pain when used correctly and for a limited time. Most people take it briefly after surgery or injury without long-term problems. The risks rise with prolonged use, and chronic pain is better managed with a combination of non-drug treatment, non-opioid medicines, and psychological support. Using dihydrocodeine under medical supervision and for the shortest time needed gives the best balance of benefit and risk.

\section*{Prevention and Self-Care}
\begin{itemize}
    \item Take dihydrocodeine exactly as prescribed
    \item Do not increase the dose without medical advice
    \item Do not drink alcohol while taking it
    \item Do not drive or operate machinery while affected
    \item Manage constipation with fluids, fibre, and laxatives as advised
    \item Store it securely, out of reach of others
    \item Return unused tablets to a pharmacy for safe disposal
    \item Discuss a plan for stopping safely with your doctor
\end{itemize}

\section*{Sources and Further Reading}
The following reputable sources provide further information for healthcare students and patients seeking reliable, current advice about dihydrocodeine.

\begin{itemize}
    \item Healthdirect Australia. \textit{Dihydrocodeine: information.} https://www.healthdirect.gov.au/
    \item NPS MedicineWise. \textit{Opioid medicines: safe use.} https://www.nps.org.au/
    \item Therapeutic Goods Administration. \textit{Opioid safety information.} https://www.tga.gov.au/
    \item NHS. \textit{Codeine and dihydrocodeine: information.} https://www.nhs.uk/medicines/
    \item MedlinePlus. \textit{Dihydrocodeine.} https://medlineplus.gov/
    \item Alcohol and Drug Foundation. \textit{Opioids: information and support.} https://adf.org.au/
\end{itemize}
